\section{The precise residual Galois realization gate}
\label{sec:s5-residual-realization}

The nonregular Galois-representation route has an exact elementary
scope test.  It supplies no new exclusion of the shared real zero
in MA-RB-002/008.
Let \(A=L_f(\,\cdot\,,\varpi)\), and let \(W\) be the specified
five-dimensional Artin representation.  Fix a finite set
\(S\) containing all bad primes of both objects.  At \(p\notin S\),
write
\[
 H_{A,p}(T)=\det(1-T\operatorname{Sat}_p\varpi),
 \qquad H_{W,p}(T)=\det(1-TW(\operatorname{Frob}_p)).
 \tag{MA-RB.107}
\]
These are the characteristic polynomials whose reciprocals
are the local Euler factors.

For a rational prime \(\ell\), let \((G_\ell)\) assert the
existence of a continuous semisimple representation
\[
 r_\ell:G_{\mathbb Q}\longrightarrow
                \operatorname{GL}_5(\overline{\mathbb F}_\ell)
 \tag{MA-RB.108}
\]
with the discrete topology on the target, unramified outside
some finite set \(S_\ell\supset S\cup\{\ell\}\), and satisfying
\[
 \det(1-Tr_\ell(\operatorname{Frob}_p))
                 =H_{A,p}(T)\pmod\ell
                 \quad(p\notin S_\ell).
 \tag{MA-RB.109}
\]
The polynomial on the right is integral, as verified below.
The finite set in this definition may depend on \(\ell\).

\paragraph{MA-RB-010 (odd residual realization is an equivalent gate).}
For every fixed odd prime \(\ell\), the following conditions
are equivalent:
\begin{enumerate}
\item the specified all-place comparison \(A=L_W\) holds;
\item \(\Omega\) is finite, or equivalently the good-prime
      polynomials in (MA-RB.107) agree outside a finite set;
\item \((G_\ell)\) holds.
\end{enumerate}
By contrast, \((G_2)\) holds without the comparison hypothesis.
Thus constructing a compatible representation in any one odd
residual characteristic would settle the comparison, but its
existence is an equivalent formulation of the remaining gate.
The result uses neither \(\mathsf A\), \(M\), nor
\((\mathcal R_B)\), and does not prove \((\mathcal R_B)\).

\begin{proof}
The exact good-prime alternative is the already audited
Calegari theorem following his similitude calculation:
\(H_{A,p}=H_{W,p}\) outside the Dirichlet-density-zero set
\(\Omega\), whereas at \(p\in\Omega\) the eigenvalue
multisets are respectively
\(\{1,i,-i,1,1\}\) and \(\{1,i,-i,-1,-1\}\).
Consequently
\[
 \begin{aligned}
 H_{A,p}(T)&=(1-T)^3(1+T^2),\\
 H_{W,p}(T)&=(1-T)(1+T)^2(1+T^2),\\
 H_{A,p}(T)-H_{W,p}(T)&=-4T(1-T)(1+T^2)
                         \quad(p\in\Omega).
 \end{aligned}
 \tag{MA-RB.110}
\]
In particular the two polynomials are distinct in every odd
characteristic and equal in characteristic two.
At other good primes the Artin polynomial is integral: its
roots are roots of unity, while the rational character of
\(W\), evaluated on all powers of Frobenius, makes its
coefficients rational by Newton identities.  Rational
algebraic integers are integers.  Formula (MA-RB.110) proves
integrality at the remaining good primes as well.

Fix a number-field realization of the finite-image
representation \(W\), and a prime above \(\ell\).
A stable integral lattice is obtained by summing the finitely
many translates of an initial lattice and localizing at that
prime.  Reduction gives a continuous finite-image
representation \(\overline W_\ell\); semisimplification does
not change any characteristic polynomial.  Indeed a
composition series makes every group element block upper
triangular, so its characteristic polynomial is the product
of those on the successive quotients.

Suppose \((G_\ell)\) holds.  Both \(r_\ell\) and
\(\overline W_\ell\) have finite image: a continuous image
of a profinite compact group in a discrete space is finite.
Their product therefore factors through a common finite
Galois quotient \(G\) of \(G_{\mathbb Q}\).
Let \(D\subseteq G\) be the set on which their
characteristic polynomials differ.  It is a union of
conjugacy classes.  If \(D\ne\varnothing\), Chebotarev's
theorem gives the unramified rational primes whose Frobenius
lies in \(D\) the positive Dirichlet density
\(\lvert D\rvert/\lvert G\rvert\).
Outside the finite excluded set, however, each such prime
belongs to \(\Omega\): at every other prime both residual
polynomials equal the reduction of \(H_{W,p}\).
This contradicts the zero Dirichlet density of \(\Omega\).
Hence \(D=\varnothing\).

It follows that the two polynomials in (MA-RB.107) are
congruent at every common good prime.  When \(\ell\) is odd,
(MA-RB.110) excludes every prime of \(\Omega\) outside the
finite excluded set.  Therefore \(\Omega\) is finite.
Conversely, if \(\Omega\) is finite, the semisimplification
of \(\overline W_\ell\) satisfies (MA-RB.109) after adding
\(\Omega\) and its finite ramification set to \(S_\ell\).
This proves the equivalence of the second and third
conditions.  The equivalence of the first and second is the
previously checked finite-exception comparison, Calegari's
Lemma~4.10, already used at the start of MA-RB-002; no
new all-place matching premise is being added here.

Finally (MA-RB.110) vanishes modulo two.  Thus
\(\overline W_2^{\mathrm{ss}}\) satisfies (MA-RB.109)
for all common good primes whether or not \(\Omega\)
is finite.  This proves the assertion about \((G_2)\).
\end{proof}

\paragraph{What the primary construction theorems supply.}
The source comparison is exact enough to check the numerical
Hecke hypotheses without building a new Galois representation.
Delete the fixed finite set \(S\), and let \(N^*\) be an integer
whose prime divisors contain \(S\).  For this retained
Hecke-polynomial collection the coefficient field is
\(\mathbb Q\), the coefficients are integral, and every
field-conjugate collection is the original one.
If \(a_j(p)\) is its coefficient of degree \(j\), every
Satake root has modulus one, so
\[
 |a_j(p)|\le \binom{5}{j},\qquad
 \sum_{p\nmid N^*}|a_j(p)|^2p^{-\sigma}
 \le \binom{5}{j}^{\!2}
                   \log\frac1{\sigma-1}+O(1)
                         \quad(\sigma\downarrow1).
 \tag{MA-RB.111}
\]
For the last bound, subtract \(\sum_p p^{-\sigma}\)
from the Euler logarithm of \(\zeta(\sigma)\).
The sum over powers at least two is bounded uniformly
for \(\sigma>1\), and
\(\log\zeta(\sigma)=\log(1/(\sigma-1))+O(1)\).
Removing finitely many primes changes only the bounded term.

These are the finite-set versions of assumptions (1), (2),
(3), and (5) in Kim and Yamauchi's \(\operatorname{GL}_n\)
construction.  Their Conjecture~5.1 and assumption (4) require
actual compatible semisimple residual Galois representations.
Section~6 reconstructs Artin representations after that input;
it does not construct the input.  The finite-set formulation
is explicit here: no rationality assertion at an omitted place
is hidden in (MA-RB.111), and no application of their complete
theorem without its Galois hypothesis is claimed.
MA-RB-010 shows that even one odd-characteristic instance of
the needed compatibility already contains the comparison
obstruction in this family.

The earlier Kim--Yamauchi \(\operatorname{GSp}_4\) theorem
also assumes residual representations, together with transfer,
integrality, and rationality of the real-analytic Hecke module.
Its stated infinity type is not a holomorphic limit of
discrete series.  Calegari--Geraghty's displayed nonregular
modularity-lifting theorems instead concern holomorphic
weight \((j,2)\), with \(j\ge4\), and Hodge--Tate weights
\([0,0,j-1,j-1]\), besides ordinary and residual-image
conditions.  These are not the all-zero weights of the
finite-image spin candidate.  Thus these precise inspected
theorems do not supply the missing representation here.
This is a hypothesis check, not a claim that no future or
other construction can do so.

\paragraph{Finite-cover interpretation and exact locators.}
The known Artin object \(W\) comes from a representation
of the finite Galois cover after the ramified primes are
removed.  Its integral lattice and residual realizations
are available independently of \(A\).  The proof uses the
common finite quotient of two such continuous residual
representations and its Frobenius classes.  It does not
assign an Artin motive to \(A\) from its finite list of
Hecke polynomials.  The inherited categorical locators are
SGA~1, Expos\'e~V, \S7, and Deligne (1979),
\emph{Valeurs de fonctions \(L\) et p\'eriodes d'int\'egrales},
Definition~6.1 and \S\S6.2--6.3; they are not being cited
as a construction of the missing compatible realization.

Calegari (2013), Theorem~4.7 and Lemma~4.10, supplies the
local alternative, density, and finite-exception closure;
the exact supplied LaTeX spans are
\texttt{S5ArtinWeb.tex}, lines 736--765 and 823--854.
Kim and Yamauchi (2016),
\emph{A conditional construction of Artin representations
for real analytic Siegel cusp forms of weight \((2,1)\)},
Theorem~1.1 and \S3.4, are in arXiv:1307.7619v5,
pp.~2--4 and 14--15,
\url{https://arxiv.org/pdf/1307.7619}; chapter DOI
\url{https://doi.org/10.1090/conm/664/13061}.
Their (2017) \emph{A uniform structure on subgroups of
\(\operatorname{GL}_n(\mathbb F_q)\) and its application to a
conditional construction of Artin representations of
\(\operatorname{GL}_n\)}, \emph{Journal of the Ramanujan
Mathematical Society}, 32(1), 75--99, is arXiv:1403.5535v3:
Theorem~1.1, pp.~2--3; Conjecture~5.1 and \S6,
pp.~13--16; and \S\S7--8, pp.~16--18,
\url{https://arxiv.org/pdf/1403.5535}.
Calegari and Geraghty (2020), \emph{Minimal modularity
lifting for nonregular symplectic representations},
\emph{Duke Mathematical Journal, 169}(5), 801--896,
Theorems~1.1--1.2, published pp.~802--803
(author-hosted advance PDF pp.~2--3),
\url{https://www.math.uchicago.edu/~fcale/papers/Siegel.pdf},
DOI \url{https://doi.org/10.1215/00127094-2019-0044},
supplies the different-weight hypothesis comparison.
The advance copy carries a 2019 copyright notice; the
published article is dated 2020.
The scope statements above do not certify complete independent
proof checks of those construction or modularity theorems.
